\section{Validation}
\label{sec:validation}

We validate the three theorems on simulated data with known
instance-level ground truth. The simulation is implemented in
\texttt{scripts/sim.R} and \texttt{scripts/run.R}, fully reproducible
from seeds \texttt{20260523} through \texttt{20260542} (twenty
seeds). Every headline number below is reported as a median with
inter-quartile range across seeds, so the reader can tell which
quantities are stable point estimates and which are
seed-dependent. Real-data application to MUSK1 and MUSK2 follows in
\cref{sec:musk}.

\subsection{Data-generating process}

For $K$ instance types and $N$ bags, per-type positivity $\eta_k$ is
fixed, bag size $|B_i| \sim \mathrm{Poisson}(\lambda)$, and the
composition $\bm{m}_i$ is drawn by assigning each instance a type
uniformly at random. Conditional on $\bm{m}_i$, latent instance
labels $y_{ij} \sim \mathrm{Bernoulli}(\eta_{k(j)})$, and the bag
label is set by the chosen aggregation rule (deterministic OR,
noisy-OR, label-noise, or $r$-of-$n$). The estimator is the
bag-OR maximum-likelihood fit of \cref{sec:identifiability},
implemented in R as a box-constrained L-BFGS-B problem in
$\bm{s} \geq 0$ with the analytic gradient
$\partial_{s_k}\log L = \sum_i m_{ik}(Y_i - p_i)/p_i$ from
\cref{thm:bag-total}.

\subsection{Rank condition (\cref{thm:rank})}

We compare three designs at $K = 5$,
$\bm{\eta} = (0.05, 0.10, 0.20, 0.30, 0.50)$, $N = 600$ bags.

\paragraph{Full-rank design.} Bag compositions are drawn uniformly
across types. The composition matrix has $\mathrm{rank}(M) = 5$ on
every seed, and the per-type recovery error has cross-seed median
$\max_k|\hat\eta_k - \eta_k| = 0.088$ with IQR $[0.066, 0.107]$:
the sampling-noise floor at this $N$.

\paragraph{Rank-deficient design.} Types 1 and 2 are forced into the
fixed co-occurrence $m_{i1} = 2 m_{i2}$ in every bag, so column 1 of
$M$ equals twice column 2; types 3--5 vary independently. The
composition matrix has $\mathrm{rank}(M) = 4$ on every seed, and
only the linear functional $2 s_1 + s_2$ is identified. Across
three optimizer starts that perturb the $(s_1, s_2)$ split,
$\hat s_1$ varies by a median $0.060$ (IQR $[0.053, 0.060]$) while
$2\hat s_1 + \hat s_2$ collapses to a spread with median
$1.6\times 10^{-15}$ (IQR $[4\times 10^{-16}, 3\times 10^{-6}]$):
machine precision in the bulk, with a small upper tail where
optimizer-tolerance jitter dominates. The recovered combination
across seeds has median $0.215$ (IQR $[0.158, 0.259]$), with truth
$2 s_1^* + s_2^* = 0.208$. This is the empirical signature of the
rank theorem: distinct $(s_1, s_2)$ pairs yield identical
bag-label distributions, but the data-identified projection onto
$\mathrm{col}(M)$ is recovered up to sampling noise.

\paragraph{Singleton-bag augmentation.} Adding $20$ singleton bags
per type to the rank-deficient design (each carrying a single
instance whose Bernoulli label is observed directly) restores
$\mathrm{rank}(M) = 5$ on every seed. The spread of $\hat s_1$
across starts collapses by more than three orders of magnitude to
median $4.8\times 10^{-5}$ (IQR $[3.5, 6.3]\times 10^{-5}$), and
the per-type recovery error returns to the sampling-noise floor,
median $0.088$ with IQR $[0.060, 0.100]$, comparable to the
full-rank baseline. \Cref{fig:rank-collapse} visualises the
collapse: under rank-deficiency, the fitted $(\hat s_1, \hat s_2)$
pairs from three starts per seed scatter along the identified-combo
ridge $2 s_1 + s_2 = 0.208$; singleton augmentation pins them to a
tight cloud around the truth.

\begin{figure}[t]
\centering
\includegraphics[width=0.95\linewidth]{fig_rank_collapse.pdf}
\caption{Rank-collapse diagnostic. Each point is one
$(\hat s_1, \hat s_2)$ from one of three optimizer starts on one of
20 seeds. \textbf{Left}: rank-deficient design ($m_{i1} = 2 m_{i2}$
in every bag); fits scatter along the dashed identified ridge
$2 s_1 + s_2 = 0.208$, individual coordinates non-identified.
\textbf{Right}: same design augmented with 20 singleton bags per
type; fits collapse to a tight cluster around the truth (cross).}
\label{fig:rank-collapse}
\end{figure}

\subsection{Bag-prevalence consistency (\cref{thm:bag-total})}

We fit the bag-OR MLE on $30$ replicates per seed at $K = 4$,
$\bm\eta = (0.05, 0.15, 0.30, 0.60)$, $N = 400$, and verify
\eqref{eq:bag-total} numerically. Boundary maximizers (with at
least one $\hat s_k$ on the lower bound) account for a median $7$
of $30$ replicates per seed (IQR $[6, 9]$), with the rest being
interior MLEs.

\paragraph{Interior MLEs.} The infinity-norm of the IRLS-weighted
residual $M^\top D^{-1}(\bm Y - \hat{\bm p})$ has cross-seed median
$3.6\times 10^{-7}$ (IQR $[3.0, 4.6]\times 10^{-7}$), with the
per-seed maximum at median $2.0\times 10^{-6}$ (IQR
$[1.6, 2.8]\times 10^{-6}$): at optimizer tolerance, confirming the
score equation \eqref{eq:bag-total}.

\paragraph{Contrast with the unweighted residual.} On the same
fits, the inf-norm of $M^\top(\bm Y - \hat{\bm p})$ has cross-seed
median $4.66$ (IQR $[4.32, 4.93]$), with per-seed maximum at median
$11.5$ (IQR $[10.3, 12.6]$): the unweighted moment-matching
identity is \emph{not} satisfied, by seven orders of magnitude
relative to the weighted form. This empirical contrast is the
practical reason the log-survival-link non-canonicity in the remark
following \cref{thm:bag-total} matters.

\paragraph{Boundary MLEs.} Boundary replicates have weighted
residuals of order $1$ to $30$ on the active coordinate, as the
KKT conditions allow (the inactive sign constraint absorbs the
remaining gradient). At interior maximizers the theorem holds; at
boundary maximizers \eqref{eq:bag-total} reduces to a KKT
inequality on the constrained coordinate.

\subsection{Aggregation-rule bias (\cref{thm:bias-rule})}

We sweep two of the three rules considered in \cref{sec:methodology}
at $K = 4$, $N = 800$, $25$ replicates per cell per seed.

\paragraph{Noisy-OR firing-rate confound.} Generating with firing
rate $\rho \in \{1.0, 0.9, 0.7, 0.5\}$ and fitting the
deterministic-OR model, the cross-seed median of the per-type
ratio $\hat\eta_k / \eta_k$ tracks $\rho$ closely. At $\rho = 1.0$
the four type-medians are $(0.99, 1.01, 1.01, 1.00)$; at
$\rho = 0.9$ they are $(0.83, 0.89, 0.91, 0.91)$; at $\rho = 0.7$
they are $(0.73, 0.69, 0.70, 0.70)$; and at $\rho = 0.5$ they are
$(0.52, 0.51, 0.51, 0.50)$. \Cref{fig:firing-rate} plots the same
trajectory: ratio collapses onto the identity line $y = \rho$ for
every type, with IQR ribbons across seeds visible only near the
boundary types. This is the firing-rate confound of
\cref{rem:firing-rate}: the deterministic-OR estimator recovers
$\rho\eta$, not $\eta$, when applied to noisy-OR data. The
product is identified; the factors are not.

\begin{figure}[t]
\centering
\includegraphics[width=0.62\linewidth]{fig_firing_rate.pdf}
\caption{Firing-rate confound under noisy-OR data. Lines: median
across 20 seeds of $\hat\eta_k / \eta_k$ for the deterministic-OR
MLE applied to noisy-OR data at firing rate $\rho$; ribbons: IQR
across seeds; dashed: identity $y = \rho$. The four instance types
all collapse onto the identity, confirming that
$\hat\eta \approx \rho\eta$: the deterministic-OR estimator
recovers only the product.}
\label{fig:firing-rate}
\end{figure}

\paragraph{Label noise.} Generating with symmetric flip rate
$\epsilon \in \{0, 0.05, 0.10, 0.20\}$, the cross-seed median bias
$\hat\eta_k - \eta_k$ grows monotonically with $\epsilon$, most
sharply for the highest-$\eta$ type. The four type-medians at
$\epsilon = 0$ are $(+0.001, +0.005, +0.000, -0.002)$ (sampling
noise); at $\epsilon = 0.05$, $(+0.023, +0.001, -0.033, -0.104)$;
at $\epsilon = 0.10$, $(+0.036, -0.001, -0.059, -0.187)$; and at
$\epsilon = 0.20$, $(+0.055, -0.010, -0.102, -0.297)$. For the
highest-$\eta$ type the bias is well-approximated as linear in
$\epsilon$ at $\epsilon \leq 0.10$ (the first-order Taylor
prediction of \eqref{eq:bias-rule}) and grows super-linearly at
$\epsilon = 0.20$ where the second-order term becomes substantial.

\subsection{Real-data application: MUSK1 and MUSK2}
\label{sec:musk}

We apply the framework to the founding MIL benchmark
\citep{dietterich1997solving}: $92$ molecules in MUSK1 ($476$
conformations) and $102$ molecules in MUSK2 ($6{,}598$
conformations), with $166$-dimensional chemical descriptors per
conformation and a binary musky/non-musky label per molecule.
Conformations are the instances; molecules are the bags. We apply
the discrete-type version of \cref{thm:rank} by z-scoring the $166$
descriptors and clustering instances with $k$-means at $K = 20$
(results at $K = 10$ are qualitatively similar; see
\texttt{scripts/musk.R}); the continuous-feature version of
\cref{sec:continuous} is reported alongside. The bag-by-type composition
matrix $M$ is built from cluster assignments, and the bag-OR MLE
of \cref{thm:bag-total} is fit. Code reproduces from
\texttt{scripts/musk.R} (seed \texttt{20260527}).

\paragraph{Rank-condition diagnostic.} The $K$-means composition
matrix has $\mathrm{rank}(M) = 20$ for both MUSK1 (condition number
$26$) and MUSK2 (condition number $93$). \Cref{thm:rank} therefore
licenses identifiability of all per-type positivities from bag
labels alone; the high MUSK2 condition number flags that the
problem is well-posed but numerically delicate.

\paragraph{Per-type positivity.} On MUSK1 the fitted $\hat\eta_k$
range from $0.58$ (the most musky cluster) down to a tail of types
at the lower bound $\hat s_k = 0$, with $11$ of $20$ types on the
boundary, exactly the KKT regime characterised in
\cref{thm:bag-total}. The interpretation is that those cluster IDs
appear only in negative bags, so the bag-OR likelihood pushes their
instance positivity to zero. On MUSK2 the boundary count is $17$ of
$20$, reflecting the much larger ratio of conformations to
positive molecules. \Cref{fig:musk-eta} plots the sorted
$\hat\eta_k$ for both datasets, with boundary types shown in light
grey.

\begin{figure}[t]
\centering
\includegraphics[width=0.95\linewidth]{fig_musk_eta.pdf}
\caption{MUSK per-instance-type positivity. Each bar is one
$k$-means cluster of standardized $166$-dim conformation
descriptors, sorted by fitted $\hat\eta_k$. Light bars mark
boundary types ($\hat s_k = 0$) where the bag-OR MLE collapses to
zero instance positivity. The long tail of boundary types is the
KKT-regime signature of \cref{thm:bag-total} on real data.}
\label{fig:musk-eta}
\end{figure}

\paragraph{Bag-prevalence consistency, real data.} \Cref{thm:bag-total}
is an \emph{interior}-MLE identity (the weighted score $M^\top
D^{-1}(\bm Y - \hat{\bm p}) = \bm 0$), so it constrains the fit only
where coordinates of $\hat{\bm s}$ are interior. On MUSK the fit is
dominated by the boundary: only $9$ of $20$ instance types on MUSK1
and $3$ of $20$ on MUSK2 are interior at the optimum, and
accordingly the weighted score residual is large
($\|M^\top D^{-1}(\bm Y - \hat{\bm p})\|_\infty = 26$ on MUSK1 and
$612$ on MUSK2), exactly because the KKT conditions, not the
interior identity, govern the constrained coordinates. The theorem
is therefore mostly \emph{inapplicable} on these heavily-boundary
fits rather than confirmed by them; the simulation of
\cref{sec:validation} is where the interior identity is verified
cleanly. As a separate and coarser observation, the overall fitted
bag-positive rate still tracks the observed rate ($\bar{\hat p} =
0.49$ vs. $\bar Y = 0.51$ on MUSK1; $0.24$ vs. $0.38$ on MUSK2),
the larger MUSK2 gap reflecting the same boundary dominance.

\paragraph{Bag-level prediction.} Under leave-one-bag-out
cross-validation (a fresh $k$-means on each held-in fold), the
bag-OR MLE achieves LOO AUC $0.79$ and accuracy $0.75$ on MUSK1, and
LOO AUC $0.60$ on MUSK2. Published continuous-feature methods
reach $\sim 0.87$ accuracy on MUSK1 \citep{andrews2003support} and
$\sim 0.84$ on MUSK2 \citep{maron1998framework}. The gap is the
cost of discretization: $k$-means on $166$-dimensional descriptors
collapses fine chemical distinctions that the discrete-type
theorem does not see. The continuous-feature version of
\cref{thm:rank} (\cref{thm:rank-continuous}), in which $M$ is
replaced by the bag-summed feature matrix $\Phi$, lifts the
discretization constraint but at $d = 166 > N$ enters the
under-identified regime; \cref{sec:musk-continuous} reports the
empirical consequence.

\paragraph{Comparison to published baselines.}
\Cref{tab:musk-comparison} places the bag-OR MLE alongside the
canonical MIL methods on MUSK. The discrete-type fit sits well
below all continuous-feature methods on MUSK1 (LOO accuracy $0.75$
vs $\geq 0.84$) and substantially below on MUSK2 (LOO AUC $0.60$).
This is the expected discretization gap: \cref{thm:rank} as stated
is for $K$ discrete types, $k$-means at $K = 20$ collapses
chemical-descriptor variation, and the discrete-type
identifiability machinery does not by itself supply a competitive
predictor on continuous features. \Cref{sec:musk-continuous}
reports the continuous-feature ($d = 166$) refit and shows that
the gap is more nuanced than discretization alone: the
identifiability diagnostic catches a second failure mode in the
$d > N$ regime.

\begin{table}[t]
\centering
\small
\begin{tabular}{lll}
\toprule
\textbf{Method} & \textbf{MUSK1} & \textbf{MUSK2} \\
\midrule
iAPR \citep{dietterich1997solving}            & 0.92 LOO acc.  & 0.89 LOO acc. \\
Diverse Density \citep{maron1998framework}    & 0.89 LOO acc.  & 0.82 LOO acc. \\
MI-SVM \citep{andrews2003support}             & 0.87 10-CV acc.& 0.84 10-CV acc. \\
mi-SVM \citep{andrews2003support}             & 0.78 10-CV acc.& 0.84 10-CV acc. \\
\midrule
Bag-OR MLE, $k$-means $K = 20$ (this paper)   & 0.75 LOO acc.  & 0.56 LOO acc. \\
                                              & 0.79 LOO AUC   & 0.60 LOO AUC \\
\bottomrule
\end{tabular}
\caption{MUSK benchmark accuracy. Published numbers are cited from
the respective papers and use the evaluation protocol each paper
reports (LOO molecule cross-validation for iAPR and Diverse
Density; 10-fold CV for MI-SVM and mi-SVM). Bag-OR MLE results in
the last row use leave-one-bag-out CV with a fresh $k$-means on
each held-in fold. The gap is consistent with the
discretization cost discussed in the prediction paragraph.}
\label{tab:musk-comparison}
\end{table}

\paragraph{What the discrete-type application demonstrates.} The
discrete-K results are not a competitive prediction benchmark.
They show that (i) the rank-condition diagnostic computes cleanly
on real chemical-descriptor data and rules out co-occurrence
confound; (ii) real chemical-descriptor data drive the bag-OR MLE
into the boundary regime (most instance types inactive at the
optimum), so \cref{thm:bag-total}'s interior identity is largely
inapplicable here and the clean confirmation lives in the
simulation, a distinction the framework makes explicit;
(iii) the interior-versus-boundary MLE
classification from \cref{sec:validation}'s simulation survives the
transition to real data; and (iv) the discrete-type framework's
predictive ceiling on continuous-feature chemistry data is real
and identifies the continuous-feature extension as the principled
next step rather than a parameter-tuning exercise.

\subsection{MUSK refit with continuous-feature \cref{thm:rank-continuous}}
\label{sec:musk-continuous}

We now refit MUSK with the continuous-feature version of the
framework, applying \cref{thm:rank-continuous} to the same data.
Each conformation's $166$-dim descriptor is min-max scaled per
column to $[0, 1]$, yielding $\bm\phi(x) \in [0, 1]^{166}$, and the
bag-summed feature matrix $\Phi = [\bm\Phi_1 \mid \cdots \mid
\bm\Phi_N]^\top \in \R^{N \times 166}_{\geq 0}$ is the
continuous-feature analogue of the discrete composition matrix
$M$.

\paragraph{Rank diagnostic (continuous).} For both MUSK1 ($N=92$)
and MUSK2 ($N=102$), $d=166 > N$, so $\mathrm{rank}(\Phi) \leq N$
by construction; the continuous parameter $\bm\beta$ is
under-identified by at least $d - N$ dimensions
(\cref{thm:rank-continuous}). Bag-level predictions, which depend
on $\bm\beta$ only through the linear predictor $\Phi\bm\beta$,
remain well-defined on the unidentified null space, and so does the
LOO-CV evaluation.

\paragraph{Bag-level prediction (continuous).} Under leave-one-bag-
out CV with a fresh per-fold min-max scaling, the continuous-feature
bag-OR MLE achieves the following on the same fold structure as
\cref{sec:musk}'s discrete-K analysis:
\begin{center}\small
\begin{tabular}{lll}
\toprule
                                                   & MUSK1                          & MUSK2 \\
\midrule
Discrete-type ($K = 20$)                           & \musonemkacc{} acc / \musonemkauc{} AUC  & \musttwomkacc{} acc / \musttwomkauc{} AUC \\
Continuous-feature ($d = 166$)                     & \musonemcacc{} acc / \musonemcauc{} AUC  & \musttwomcacc{} acc / \musttwomcauc{} AUC \\
Continuous-feature, PCA-to-$k=50$ (\cref{sec:musk-pca}) & \musonempcacc{} acc / \musonempcauc{} AUC & \musttwompcacc{} acc / \musttwompcauc{} AUC \\
\bottomrule
\end{tabular}
\end{center}
The continuous-feature refit \emph{\musresumphrase{}} the
discrete-$K=20$ fit on both datasets. This is the empirical
signature of the under-identification \cref{thm:rank-continuous}
predicts when $d > N$: each LOO refit lands at a different point on
the $74$-dimensional (MUSK1) or $64$-dimensional (MUSK2) null space
of $\Phi$, and held-out bag predictions inherit the variation
because the feature scaling, the held-out instances, and the chosen
null-space point all shift fold-to-fold. The discrete-$K = 20$ fit
throws away chemical-descriptor variation but has an identified
parameter on every fold, so its LOO predictions are consistent.

\paragraph{Interpretation.} The framework's value on real data is
twofold. First, the rank diagnostic remains a precise structural
tool under either feature map; it correctly warned that the
$d = 166$ continuous-feature fit on $N \leq 102$ bags would be
under-identified, and the LOO degradation confirmed the warning.
Second, the framework's own prescription for repair is
rank-reduction of the feature map so that $\mathrm{rank}(\Phi)$
equals the feature dimension. \Cref{sec:musk-pca} carries out that
prescription on MUSK with PCA to $k = 50$ principal components.

\subsection{PCA-to-$k$ rank-reduction (framework prescription)}
\label{sec:musk-pca}

To test the framework's prescription empirically, we project the
min-max-scaled features onto the top $k = 50$ principal components
(refit per fold) and shift each component to be nonnegative so the
$\bm\beta \geq 0$ constraint stays well-posed. The bag-summed
feature matrix $\Phi_{\mathrm{PCA}} \in \R^{N \times 50}_{\geq 0}$
then has rank $50$ for both MUSK1 and MUSK2 ($k < N$ in both
cases), so \cref{thm:rank-continuous} guarantees identifiability of
$\bm\beta$.

\paragraph{Bag-level prediction (PCA).} The PCA-reduced
continuous-feature MLE achieves LOO accuracy $\musonempcacc{}$ and
LOO AUC $\musonempcauc{}$ on MUSK1, and accuracy
$\musttwompcacc{}$ and AUC $\musttwompcauc{}$ on MUSK2. Relative to
the under-identified $d = 166$ fit of the previous subsection, the
identified PCA-to-$k = 50$ fit improves on both datasets (MUSK1 AUC
$0.66 \to 0.68$, MUSK2 AUC $0.54 \to 0.56$), in the direction
\cref{thm:rank-continuous} predicts: restoring identifiability
stabilizes the LOO refits. The PCA fit does not, however, reach the
discrete-$K = 20$ baseline of \cref{sec:musk} (MUSK1 AUC $0.79$,
MUSK2 AUC $0.60$), and remains well below the published
continuous-feature methods of \cref{tab:musk-comparison}.

\paragraph{What the PCA experiment confirms, and what it does not.}
The identifiability chain holds end to end: (i) the diagnostic
flagged the $d > N$ regime as under-identified; (ii) the LOO of the
$d = 166$ fit degraded as predicted; (iii) rank-reduction restored
an identifiable parameter, $\mathrm{rank}(\Phi_{\mathrm{PCA}}) = 50
= d$, as \cref{thm:rank-continuous} requires; and (iv) the LOO of
the identified PCA fit improved over the under-identified fit. The
framework's diagnostic and prescriptive content, that
under-identification hurts predictive stability and that
rank-reduction repairs it, is confirmed on real data.

What the experiment does \emph{not} show is that identifiability
buys competitiveness. The PCA fit stays below the discrete-$K = 20$
baseline and far below mi-SVM and Diverse Density. Identifiability is
necessary for stable estimation but not sufficient for predictive
performance: the bag-OR likelihood with a linear-in-principal-
components hazard is a weaker model class than the max-margin and
concept-point methods, independent of whether its parameter is
identified. The framework's contribution is to separate these two
axes cleanly. It tells the practitioner when a parameter can be
estimated at all (the rank condition) and when an estimate is
stable (identifiability plus $\mathrm{rank}(\Phi) = d$), and it
leaves the choice of hazard model class as a distinct decision that
the identifiability analysis deliberately does not prejudge.

\subsection{Future-work benchmarks}

Elephant/Fox/Tiger image-region MIL \citep{andrews2003support} and
weakly supervised computational pathology
\citep{campanella2019clinical,lu2021data,bejnordi2017diagnostic} are
natural next applications of the framework, the latter with expert
pixel annotations serving as singleton candidate sets for the
firing-rate confound of \cref{rem:firing-rate}. Both are left to
future work; the MUSK demonstration in \cref{sec:musk} is
sufficient to establish the structural diagnostics on real data.
